\chapter{Results and Evaluation}

This chapter reports the measured behaviour of the system across four
dimensions: the wireless link, on-device model performance, recognition
accuracy, and power/thermal robustness. Where a result has an important
limitation, it is stated plainly.

\section{Wireless Link Reliability}

With the master in BLE mode and all four slaves powered, the master's periodic
diagnostics report the number of packets received from each finger over each
five-second window. In a healthy run every finger reports approximately 125
packets per five seconds—exactly the expected $25\,\text{Hz} \times 5\,\text{s}$
—with the per-finger loss counters at zero (Table~\ref{tab:rf}). The BLE link,
once the MTU is negotiated to 256 bytes, delivers the full 180-byte feature
vector per notification at the same 25\,Hz. The architecture change from a
wrist hub with external antennas to nail-mounted modules on the on-chip
antennas was decisive here: at the few-centimetre finger-to-finger distance the
link is no longer pose-dependent, and packet loss in the closed-fist
configuration—previously the worst case—is negligible.

\begin{table}[!htbp]
\centering
\caption{\label{tab:rf}Representative five-second link diagnostics from the
master (BLE connected).}
\begin{tabular}{lcc}
\toprule
Finger & Packets received & Lost \\
\midrule
Thumb  & 125 & 0 \\
Index  & 125 & 0 \\
Middle (own) & 125 & 0 \\
Ring   & 125 & 0 \\
Pinky  & 125 & 0 \\
\bottomrule
\end{tabular}
\end{table}

\section{On-Device Model Performance}

The deployed int8 model is small and fast (Table~\ref{tab:perf}). At roughly
1\,ms per inference against a 40\,ms (25\,Hz) frame budget, inference latency is
negligible, and the memory footprint is small enough for any modern phone. The
gated-capture design means inference runs once per user request rather than
continuously, so the practical compute and energy cost is lower still.

\begin{table}[!htbp]
\centering
\caption{\label{tab:perf}On-device characteristics of the quantised model
(Edge Impulse estimate, Raspberry Pi~4 class target).}
\begin{tabular}{lr}
\toprule
Metric & Value \\
\midrule
Inference time   & $\approx$ 1\,ms \\
Peak RAM         & 22.7\,KB \\
Flash (model)    & 137\,KB \\
Quantisation     & int8 (input and output) \\
\bottomrule
\end{tabular}
\end{table}

\section{Recognition Accuracy}

Several models were trained over the course of the project as the feature
representation evolved. Table~\ref{tab:acc} summarises the held-out test
accuracy of the main variants. The test sets are split off \emph{before}
augmentation, so augmented copies of test samples never appear in training and
the figures are not inflated by that form of leakage.

\begin{table}[!htbp]
\centering
\caption{\label{tab:acc}Held-out test accuracy of the main model variants.
All figures are within-session (training and test recorded in the same
session); see the discussion on cross-session generalisation.}
\begin{tabular}{lccc}
\toprule
Model & Classes & Features & Held-out accuracy \\
\midrule
Smoke set        & 5  & raw / static     & high (pipeline validation) \\
Letters v2       & 26 & relative, static & $\approx$100\% \\
Words            & 15 & absolute, dynamic& $\approx$100\% \\
Letters v3       & 26 & absolute, dynamic& high \\
\bottomrule
\end{tabular}
\end{table}

\subsection{The Single-Session Caveat}

The held-out accuracies in Table~\ref{tab:acc} are high, but they are
\emph{within-session}: each test set shares the same calibration, mounting, and
thermal state as its training set. This is the easiest possible test and must
not be read as the system's accuracy for a fresh user session. The honest
measure of generalisation is cross-session—training on one or more sessions and
testing on a different day with a fresh calibration—and the live testing
described next is the practical proxy for it.

\subsection{Live Findings and the Role of Feature Representation}

Live testing was the experiment that drove the project's main finding. The
static, absolute-feature letter model achieved high offline accuracy yet failed
live, confidently confusing letters because the calibration pose could not be
reproduced. Introducing the relative (orientation-invariant) transform removed
this calibration sensitivity, but the thumb-cluster letters
(\textbf{A}/\textbf{S}/\textbf{T}/\textbf{M}/\textbf{N}) remained confused
live, confirming that their separation is limited by sensor resolution rather
than by the feature transform. For the word model, the relative transform
introduced a different failure: signs distinguished only by global hand
orientation or translation (notably \textsc{you} vs \textsc{me}) collapsed,
because mean-centering deliberately discards exactly that information; switching
the word model to absolute features separated them again.

These observations converge on a single conclusion: for inertial sign
recognition on low-cost sensors, \emph{dynamic gestures with absolute features}
are the robust choice, because the discriminative information then resides in
the angular-rate channels, which are inherently independent of the start-up
calibration. Re-recording the alphabet as distinctive per-letter motions and
training on absolute features is the embodiment of this conclusion. A secondary
but practical finding is that a dynamic gesture must fill the capture window: a
sign with a long still tail drifts toward whichever class has the least motion,
which then acts as an attractor for ambiguous captures.

\section{Power and Thermal Behaviour}

The peak instantaneous current of the whole glove is dominated by simultaneous
radio transmissions. With the Wi-Fi transmit power at its maximum, a slave's
transmit burst draws on the order of 300\,mA; in the pathological case of all
four slaves transmitting in the same instant the system peak approaches
1.5\,A, although the phase-staggered transmissions make full alignment
momentary. The average draw is far lower, on the order of a few hundred
milliamps.

This peak is the source of the most disruptive failure observed: a marginal
supply or wiring path sags the 3.3\,V rail far enough during a transmit burst to
trip the ESP32-C3 brown-out detector and reset a module, which in the logs
appears as a repeated reboot and (on the affected board) erratic sensor reads.
Because the transmit current is the trigger, the primary mitigation is to
reduce the Wi-Fi transmit power: the finger-to-finger link spans only a few
centimetres and has tens of decibels of margin, so a much lower transmit power
both removes the current spike and reduces self-heating, with no loss of link
reliability. Improving the power wiring and adding local decoupling capacitance
at each module are complementary hardware mitigations. Reducing the transmit
power changes only the radio, not the sensor data, so it requires no
re-collection of data or re-training.

Separately, the proximity of the microcontroller to the IMU within each nail
module warms the gyroscope over a session, which shifts its zero-rate bias.
This is handled in firmware by the continuous still-state bias tracking
(Chapter~4) and is a further reason the motion-based, calibration-independent
features were preferred.
